
\chapter{Prerequisites}\label{Sec3.1}
%what is needed: 
%Cauchy Riemann equations
%fourier transforms 
%laplace transforms
%standards calculus rules and formulas you use in your main sections

This section collects the foundational results from complex analysis and real analysis that underpin the later concepts. Rather than treating these as isolated facts, the goal is to present them in the order in which they become relevant, so that each result feels motivated by what follows. 

\subsection*{Complex analytic Foundations}
\subsubsection{Cauchy-Riemann Equations}
If we write $f(z)=u(x,y)+i v(x,y)$ where $z=x+iy$, then $f$ being holomorphic is equivalent to $u$ and $v$ satisfying the Cauchy-Riemann equations: 
\begin{equation}\label{Eq3.1.2}
    \frac{\partial u}{\partial x}=\frac{\partial v}{\partial y}, \ \ \ \ \ \frac{\partial u}{\partial y}=-\frac{\partial v}{\partial x}
\end{equation} These are the conditions that tie the real and imaginary parts of $f$ together in the right way. In practice, they are most useful as a way of checking whether a given function is holomorphic, but for our purposes their real importance is theoretical — they guarantee that the Laplace transform $F(s)$, once we know it converges in a half-plane, is holomorphic there. \cite{ThousandPage, MethOfPhys, ahlfors1979}


\subsubsection*{Cauchy theorem}
If $f$ is holomorphic in a simply connected domain $D$. If $\gamma$ is any closed piecewise smooth contour lying inside $D$, then \begin{equation}\label{Eq3.1.3}
    \oint_{\gamma} f(z)dz=0
\end{equation}
What this means practically is that the integral of a holomorphic function around any closed loop is zero. For the inverse Laplace transform this is essential because it tells us we are free to deform the vertical line of integration sideways. As long as we do not cross any singularities of the integrand while doing so, the value of the integral does not change. \cite{ThousandPage,MethOfPhys,ahlfors1979}

\subsubsection*{Cauchy integral Formula}
If $f$ is holomorphic on an open set containing a simple closed contour $\Gamma$ and its interior, then for any $z_{0}$ inside $\Gamma$, 
\begin{equation}\label{Eq3.1.4}
    f(z_{0})=\frac{1}{2\pi i}\oint_{\Gamma}\frac{f(z)}{z-z_{0}}dz
\end{equation} This implies that the values of $f$ on the boundary of a region completely determines its value at every point inside. The main reason we need this here is because the residue theorem follows directly from it. \cite{YellowBook,ThousandPage,RedBook,ahlfors1979} 


\subsubsection*{Singularities and Residues}
If $f$ is holomorphic in a punctured neighbourhood $0<|z-z_{0}|<r$ but not at $z_{0}$ itself, then $z_{0}$ is called an isolated singularity. Near such a point $f$ has a Laurent expansion \begin{equation}\label{Eq3.1.5}
    f(z)=\sum^{\infty}_{n\rightarrow\infty}a_{n}(z-z_{0})^{n}
\end{equation}
Depending on what the negative power term looks like, there are three cases. 
\\ \\
A removable singularity has no negative power term at all. The function is actually bounded near $z_{0}$ and the singularity can be removed simply by defining $f(z_{0})=a_{0}$
\\ \\
A pole of order $m$ has finitely many negative power terms, down to $(z-z_{0})^{-m}$ with $a_{-m}\neq 0$. The function blows up as $z\rightarrow z_{0}$ but in a controlled way, since multiplying by $(z-z_{0})^{m}$ removes the problem entirely.
\\ \\
An essential singularity has infinitely many negative power terms. The behaviour near an essential singularity is wild in a precise sense. By the Casorati–Weierstrass theorem, the function gets arbitrarily close to every complex number in any neighbourhood of $z_{0}$.
\\ \\
The residue of $f$ at $z_{0}$ is defined as $a_{-1}$, the coefficient of the $(z-z_{0})^{-1}$ term. The reason this particular coefficient matter is that it is the only term in the Laurent series that survives integration around a closed loop: 
\begin{equation}\label{Eq3.1.6}
    \oint_{|z-z_{0}|=r}(z-z_{0})^{n}dz=\left\{ \begin{array}{lr}
    2\pi i & n=-1 \\
    0 & n\neq -1
\end{array}\right.
\end{equation}
So when we integrate around a singularity, the residue will be the only thing that contributes. \cite{RedBook, YellowBook,ahlfors1979}


\subsubsection{Estimation lemma}

If $f$ is continuous on a contour $\Gamma$ of length $L$, and $|f(z)|\leq M$ everywhere on $\Gamma$, then 
\begin{equation}\label{Eq3.1.7}
    \left|\int_{\Gamma}f(z)dz\right|\leq ML
\end{equation}
This is one of the most used tools in what follows. Whenever we need to show that an integral along some arc vanishes as the arc grows large or shrinks to a point, this is how we do it. We bound the integrand by $M$, note the arc length is $L$, and show that $ML\rightarrow0$. \cite{ahlfors1979,priestley2003}

\subsubsection*{Residue Theorem}

If $f$ is meromorphic inside and on a positively oriented closed contour $\Gamma$, with isolated poles at $z_{1},...,z_{n}$ inside $\Gamma$, then
\begin{equation}\label{Eq3.1.8}
    \oint_{\Gamma}f(z)dz= 2\pi i\sum^{n}_{k=1}\operatorname{Res}(f,z_{k})
\end{equation} For a simple pole at $z_{k}$: 
\begin{equation}\label{Eq3.1.9}
    \operatorname{Res}(f,z_{k})=\lim_{z\rightarrow z_{k}}(z-z_{k})f(z)
\end{equation} 
When $f(z)=G(z)/H(z)$ with $H(z_{k})=0$ and $H'(z_{k})\neq0$ this simplifies to: 
\begin{equation}\label{Eq3.1.10}
    \operatorname{Res}(f,z_{k})=\frac{G(z_{k})}{H'(z_{k})}
\end{equation}
For poles of order $m$: 
\begin{equation}\label{Eq3.1.11}
    \operatorname{Res}(f,z_{k})=\frac{1}{(m-1)!}\lim_{z\rightarrow z_{k}}\frac{d^{m-1}}{dz^{m-1}}[(z-z_{k})^{m}f(z)]
\end{equation}
This is the theorem that does the actual work in the complex inversion formula. Once we have closed the Bromwich contour, evaluating the inversion integral becomes a matter of computing residues. \cite{RedBook,YellowBook,WatsonLaplace,MethOfPhys,ThousandPage,TransformsTheory,ahlfors1979,priestley2003}

\subsubsection*{Jordan's Lemma (Exponential Decay)}

If $f$ is analytic in the left half-plane with $|f(s)|\leq M/|s|$ on large arcs, then for $t>0$: 
\begin{equation}\label{Eq3.1.12}
    \int_{C_{R}}e^{st}f(s)ds\rightarrow 0 \ \ \ \ \text{ as } R\rightarrow\infty
\end{equation}
Where $C_{R}$ is the left semicircular arc of radius $R$.
\\ 
The reason this works is that for $t>0$, the factor $|e^{st}|=e^{t \Re(s)}$ decays exponentially as $\Re(s)\rightarrow -\infty$ in the left half-plane. This exponential decay is strong enough to kill the contribution from the arc even as it grows large. It is worth noting that this only works for $t>0$. For $t\leq 0$ the argument breaks down entirely, which is consistent with $f(t)=0$ for $t<0$ in the inversion formula. \cite{MethOfPhys,RedBook,TransformsTheory,YellowBook}

